\documentclass[12pt,a4paper]{article}
% Estilo compartido para informes. Compilar con LuaLaTeX desde tp1/.
\usepackage[spanish,es-nodecimaldot]{babel}
\usepackage{fontspec}
\setmainfont{Montserrat-Medium.ttf}[
  Path=../assets/fonts/,
  BoldFont=Montserrat-Bold.ttf,
  ItalicFont=Montserrat-MediumItalic.ttf]
\usepackage[a4paper,margin=2.5cm,headheight=16pt]{geometry}
\usepackage{xcolor,graphicx,booktabs,tabularx}
\usepackage{float}
\usepackage{titlesec,fancyhdr,enumitem,caption}
\usepackage{ragged2e}
\usepackage{hyperref}
\definecolor{verde}{HTML}{1D4036}
\definecolor{bronce}{HTML}{D38129}
\definecolor{rosa}{HTML}{D88E77}
\definecolor{calido}{HTML}{EFE4D8}
\definecolor{grafito}{HTML}{1D1D1D}
\hypersetup{colorlinks=true,linkcolor=verde,urlcolor=verde,
  pdftitle={TP1 - Instalación de máquinas virtuales},pdfsubject={Cyberseguridad}}
\AtBeginDocument{\color{grafito}\RaggedRight\fontsize{12}{14.4}\selectfont}
\setlength{\parindent}{0pt}
\setlength{\parskip}{7pt}
\setlist[itemize]{leftmargin=*,itemsep=7pt,topsep=4pt}
\setlist[enumerate]{leftmargin=*,itemsep=6pt}
\titleformat{\section}{\color{verde}\bfseries\fontsize{18}{21.6}\selectfont}{\thesection}{0.65em}{}
\titleformat{\subsection}{\color{verde}\bfseries\large}{\thesubsection}{0.65em}{}
\titleformat{\subsubsection}{\color{verde}\bfseries}{\thesubsubsection}{0.65em}{}
\titlespacing*{\section}{0pt}{16pt}{8pt}
\titlespacing*{\subsection}{0pt}{12pt}{5pt}
\titlespacing*{\subsubsection}{0pt}{10pt}{4pt}
\setcounter{tocdepth}{2}
\pagestyle{fancy}
\fancyhf{}
\fancyhead[L]{\small\color{verde}FCEFyN | UNC}
\fancyhead[R]{\small\color{verde}\materia\ · TP \numeroTP}
\fancyfoot[L]{\scriptsize Informe de laboratorio}
\fancyfoot[R]{\small\thepage}
\renewcommand{\headrulewidth}{0.4pt}
\captionsetup{font=small,labelfont=bf,justification=raggedright,singlelinecheck=false,hypcap=false}
\newcommand{\pendiente}[1]{\par{\small\color{verde}\textbf{Por completar.} #1\par}}
% Si falta una captura se muestra un espacio reservado; no impide compilar.
\newcommand{\captura}[2]{%
  \par\smallskip\begingroup\centering
  \IfFileExists{figuras/#1}{%
    \includegraphics[width=\linewidth,height=5cm,keepaspectratio]{figuras/#1}%
  }{%
    \setlength{\fboxsep}{10pt}%
    \fcolorbox{bronce}{calido}{\parbox[c][1.25cm][c]{\dimexpr\linewidth-2\fboxsep-2\fboxrule\relax}{%
      \centering\small Captura pendiente\par\texttt{\detokenize{#1}}}}%
  }%
  \captionof{figure}{#2}\par\endgroup\smallskip
}

\usepackage{xurl}

\newcommand{\materia}{Cyberseguridad}
\newcommand{\numeroTP}{18}
\newcommand{\tituloTP}{Análisis de un infostealer (Lumma)}
\newcommand{\estudiante}{Brigido Tomas Andres}
\newcommand{\legajo}{38481096}
\newcommand{\carrera}{ICOMP}
\newcommand{\docente}{Javier A. Jorge}


% Metadatos específicos del TP18 y capturas en img/.
\hypersetup{pdftitle={TP18 - Análisis de un infostealer (Lumma)},linkcolor=black,urlcolor=black,citecolor=black}
\renewcommand{\captura}[2]{%
  \begin{figure}[H]
    \centering
    \IfFileExists{img/#1}{%
      \includegraphics[width=\linewidth,height=7cm,keepaspectratio]{img/#1}%
    }{%
      \setlength{\fboxsep}{10pt}%
      \fcolorbox{bronce}{calido}{\parbox[c][1.1cm][c]{0.85\linewidth}{%
        \centering\small Captura pendiente\par\texttt{\detokenize{#1}}}}%
    }%
    \caption{#2}
  \end{figure}
}

\usepackage{caption}
\captionsetup{justification=centering}
\begin{document}
\hypersetup{pageanchor=false}
\begin{titlepage}
  \includegraphics[width=8cm]{../assets/logo-fcefyn-2026.png}
  \vspace{1cm}

  {\color{verde}\small UNIVERSIDAD NACIONAL DE CÓRDOBA\par}
  \vspace{2.1cm}
  {\color{bronce}\bfseries\large TRABAJO PRÁCTICO \numeroTP\par}
  \vspace{0.45cm}
  {\color{verde}\bfseries\fontsize{32}{38.4}\selectfont Análisis de un\\infostealer (Lumma)\par}
  \vspace{0.3cm}
{\small PCAP, Security Onion e indicadores de compromiso\par}
\vspace{0.35cm}
  {\large \materia\par}
  \vspace{0.8cm}
  {\color{bronce}\rule{2.7cm}{2pt}\par}
  \vfill
  \renewcommand{\arraystretch}{1.45}
  \begin{tabular}{@{}p{3.2cm}p{10cm}@{}}
    Estudiante & \estudiante \\
    Legajo & \legajo \\
    Carrera & \carrera \\
    Docente & \docente \\
  \end{tabular}
  \vfill
  {\small Córdoba, Argentina\par}
\end{titlepage}
\hypersetup{pageanchor=true}

\tableofcontents
\vspace{1cm}
\newpage


\section{Introducción}
Se analizó el PCAP del caso Lumma del 19 de septiembre de 2024 para identificar comunicaciones sospechosas y cinco indicadores de compromiso. Se utilizó la VM disponible de Security Onion, con Snort, Zeek y Kibana.

\section{Proceso y resultados}
\subsection{Importación y revisión de alertas}
Se descargó el PCAP de Malware Traffic Analysis y se importó desde la carpeta Downloads con el siguiente comando:
\begin{quote}\small\ttfamily
sudo so-import-pcap\\
2024-09-19-file-downloader-to-Lumma-Stealer.pcap
\end{quote}
\captura{importacion.png}{Importación completada y análisis del PCAP con Snort y Zeek.}
En Kibana se seleccionó el intervalo del 19 al 20 de septiembre de 2024. Se visualizaron 147 registros y el panel NIDS mostró 0 alertas. Esto no demuestra ausencia de actividad maliciosa; se continuó con los registros HTTP y DNS.

\clearpage
\subsection{Análisis de HTTP y DNS}
El equipo \texttt{10.9.19.101} realizó tres solicitudes HTTP GET por el puerto 80: dos a \texttt{nhit66.com} y una a \texttt{back-kurwa.com}. Se observaron dos respuestas 200 y una redirección 301. No aparecieron solicitudes POST en el panel HTTP revisado.

Las rutas fueron \texttt{/test\_gate0117.php}, \texttt{/gate2.php} y \texttt{/gAySB.php}. Las tres solicitudes compartían un User-Agent que declaraba Windows NT 6.3 y Trident/7.0; este dato no permite identificar por sí solo el malware.
\captura{http-test-gate.png}{Registro HTTP con dominio, URI, IP de origen y User-Agent.}

Se observaron 24 registros DNS de tipo A: 16 con NOERROR y 8 con NXDOMAIN. Entre los dominios consultados aparece \texttt{back-kurwa.com}, también presente en HTTP. Los indicadores se contrastaron con el listado publicado para este caso [1].
\captura{dns-consultas.png}{Dominios consultados y respuestas DNS, mostrados como listas independientes.}

\clearpage
\section{Indicadores de compromiso}
Los siguientes cinco dominios aparecen en las capturas y en el listado de referencia [1]. Su función se atribuye al análisis publicado; no se infiere únicamente del nombre o de una consulta DNS.

\begingroup
\small
\renewcommand{\arraystretch}{1.6}
\begin{tabularx}{\linewidth}{@{}c p{5.1cm} X@{}}
\toprule
N.º & Dominio & Evidencia y relación con el caso \\
\midrule
1 & \nolinkurl{nhit66.com} & Dos GET en HTTP; tráfico inicial de la cadena de infección. \\
2 & \nolinkurl{back-kurwa.com} & Presente en HTTP y DNS; GET con redirección 301 a HTTPS según [1]. \\
3 & \nolinkurl{genedjestytw.shop} & Consultado por DNS; servidor C2 de Lumma identificado en [1]. \\
4 & \nolinkurl{eemmbryequo.shop} & Consultado por DNS; dominio C2 que no resolvió según [1]. \\
5 & \nolinkurl{chick-chick666.com} & Consultado por DNS; alojamiento del ejecutable de Lumma según [1]. \\
\bottomrule
\end{tabularx}
\par\endgroup

\section{Medidas de mitigación}
\begin{enumerate}[itemsep=5pt]
\item \textbf{Prevenir la infección:} descargar software de fuentes oficiales, evitar cracks y mantener el sistema y la protección antimalware actualizados.
\item \textbf{Detectar y contener:} monitorear HTTP y DNS, actualizar las reglas de detección y bloquear los indicadores confirmados del incidente.
\item \textbf{Proteger las cuentas:} ante una infección, aislar el equipo, revocar sesiones y cambiar contraseñas desde un dispositivo limpio; habilitar autenticación multifactor.
\end{enumerate}

\clearpage
\section{Conclusión}
El análisis permitió identificar cinco dominios vinculados al caso mediante registros HTTP y DNS y su contraste con información de referencia. Aunque NIDS no mostró alertas, la revisión del tráfico aportó evidencia útil para investigar la infección y proponer medidas de protección.

\section*{Referencia}
\addcontentsline{toc}{section}{Referencia}
[1] Malware Traffic Analysis, \textit{File downloader to Lumma Stealer}, 19 de septiembre de 2024. PCAP y listado de IOCs del caso:\\
\url{https://www.malware-traffic-analysis.net/2024/09/19/index.html}
\end{document}
